% !TEX TS-program = pdflatex
% !TEX encoding = UTF-8 Unicode
%================================================
% Animation flag
%================================================
% Set exactly ONE of the following two lines:
\newif\ifanimations
\animationstrue       % animations ON: bullet lists reveal progressively
% \animationsfalse    % animations OFF: all overlay content shown at once

\ifanimations
  \documentclass[aspectratio=169,11pt,t]{beamer}
\else
  \documentclass[aspectratio=169,11pt,t,handout]{beamer}
\fi

% Based on the style of w01a_SyllabusIntro.tex
\usetheme{SimplePlus}
\usecolortheme{default}
\usepackage[utf8]{inputenc}
\usepackage[T1]{fontenc}
\usepackage{lmodern}
\usepackage{booktabs}
\usepackage{tabularx}
\usepackage{array}
\usepackage{hyperref}
\usepackage{xcolor}
\usepackage{multicol}
\usepackage{ragged2e}
\usepackage{qrcode}

\definecolor{PittBlue}{RGB}{0,53,148}
\definecolor{PittGold}{RGB}{255,184,28}
\setbeamercolor{palette primary}{bg=PittBlue,fg=white}
\setbeamercolor{palette secondary}{bg=PittBlue!85!black,fg=white}
\setbeamercolor{palette tertiary}{bg=PittGold,fg=black}
\setbeamercolor{frametitle}{bg=PittBlue,fg=white}
\setbeamercolor{title}{fg=PittBlue}
\setbeamercolor{structure}{fg=PittBlue}
\setbeamertemplate{navigation symbols}{}

% Put every frame title in a consistent bar at the very top of the slide.
\setbeamertemplate{frametitle}{%
  \nointerlineskip%
  \begin{beamercolorbox}[wd=\paperwidth,ht=2.8ex,dp=1.1ex,leftskip=0.55cm,rightskip=0.55cm]{frametitle}%
    \usebeamerfont{frametitle}\insertframetitle
  \end{beamercolorbox}%
}
\setbeamertemplate{itemize items}[circle]
\setbeamertemplate{enumerate items}[default]

% Animation control:
%   \animationstrue  -> progressive overlays are enabled
%   \animationsfalse -> handout mode; all overlay content appears at once
\newcommand{\smallitem}{\setlength\itemsep{2pt}}
\newcommand{\tightitems}{\setlength\itemsep{1pt}\setlength\parskip{0pt}}
\newcolumntype{Y}{>{\RaggedRight\arraybackslash}X}
\newcolumntype{Z}{>{\RaggedRight\arraybackslash}p{0.30\textwidth}}

\title[Engineering Data Analytics Certificate]{Engineering Data Analytics Certificate}
\subtitle{Programming, Statistics, Machine Learning, and Engineering Applications}
\author{Prof. Paul W. Leu}
\institute{University of Pittsburgh}
\date{Fall 2026}

\begin{document}

\begin{frame}[t]
  \titlepage
\end{frame}

\begin{frame}[t]{Why Data Analytics?}
\large
Data is changing how engineers design, manufacture, optimize, and make decisions.

\vspace{0.35cm}
\begin{itemize}[<+->]
  \smallitem
  \item Companies have strong demand for engineers who can analyze data and communicate results.
  \item Data analytics skills complement every engineering major.
  \item The certificate helps document a coherent set of skills on your record.
\end{itemize}

\vspace{0.35cm}
\uncover<4->{\begin{center}
\Large \textbf{Goal:} combine engineering domain knowledge with modern data methods.
\end{center}}
\end{frame}

\begin{frame}[t]{Skills in Demand}
\begin{columns}[T,onlytextwidth]
\column{0.48\textwidth}
\textbf{Core technical skills}
\begin{itemize}[<+->]
  \tightitems
  \item Program in R or Python
  \item Analyze data with statistics
  \item Build and interpret models
  \item Use machine learning and AI tools
  \item Make clear visualizations
\end{itemize}

\column{0.48\textwidth}
\textbf{Engineering applications}
\begin{itemize}[<+->]
  \tightitems
  \item Make better decisions
  \item Optimize supply chains
  \item Manufacture products
  \item Manage energy systems
  \item Develop new applications
\end{itemize}
\end{columns}
\end{frame}

\begin{frame}[t]{Career Motivation}
\begin{columns}[T,onlytextwidth]
\column{0.49\textwidth}
\begin{block}{Projected growth}
\centering
{\Huge 34\%}\par
\vspace{0.1cm}
Projected growth in data-scientist jobs through 2034
\end{block}

\column{0.49\textwidth}
\begin{block}{Salary signal}
\centering
{\Huge \$166K}\par
\vspace{0.1cm}
Average U.S. data-scientist salary as of early 2026
\end{block}
\end{columns}

\vspace{0.45cm}
\begin{center}
\textbf{The certificate is a way to show employers that your engineering degree includes formal data-analytics training.}
\end{center}
\end{frame}

\begin{frame}[t]{Certificate Structure}
The Engineering Data Analytics Certificate requires \textbf{15 credits} organized around three domains.

\vspace{0.25cm}
\begin{center}
\begin{tabularx}{0.95\textwidth}{@{}lY@{}}
\toprule
\textbf{Domain} & \textbf{Purpose}\\
\midrule
1. Foundations & Programming and inferential statistics\\
2. Expertise & Exploratory/applied data science plus modeling and prediction\\
3. Specialization & A semester-long engineering data analytics project\\
\bottomrule
\end{tabularx}
\end{center}

\vspace{0.25cm}
\begin{itemize}[<+->]
  \item At least \textbf{9 credits} must be beyond courses required for your major.
  \item Students may count a maximum of \textbf{two major courses} toward the certificate.
\end{itemize}
\end{frame}

\begin{frame}[t]{Domain 1: Foundations}
\begin{columns}[T,onlytextwidth]
\column{0.48\textwidth}
\textbf{Data Science Programming}\par
\vspace{0.15cm}
Choose one course, such as:
\begin{itemize}[<+->]
  \tightitems
  \item CEE 0109
  \item CS 1501
  \item ECE 1148
  \item IE 0015
  \item STAT 1060
\end{itemize}

\column{0.48\textwidth}
\textbf{Inferential Statistics}\par
\vspace{0.15cm}
Choose one course, such as:
\begin{itemize}[<+->]
  \tightitems
  \item BioE 1000
  \item ECE 0402
  \item ENGR 0021
  \item IE 1071
  \item STAT 1000 / 1152 / 1361
\end{itemize}
\end{columns}
\end{frame}

\begin{frame}[t]{Domain 2: Expertise}
\begin{columns}[T,onlytextwidth]
\column{0.48\textwidth}
\textbf{Exploratory Applied Data Science}\par
\vspace{0.15cm}
Choose one course, such as:
\begin{itemize}[<+->]
  \tightitems
  \item ENGR 1451/2451: Exploratory Data Science
  \item ECE 1147 or ECE 1395
  \item INFSCI 0510
  \item MEMS 1300
  \item CS 1675
\end{itemize}

\column{0.48\textwidth}
\textbf{Modeling and Prediction}\par
\vspace{0.15cm}
Choose one course, such as:
\begin{itemize}[<+->]
  \tightitems
  \item ENGR 1453/2453
  \item IE 1187
  \item IE 1062 or IE 1072
  \item INFSCI 1530
  \item MEMS 1069 or MEMS 1145
\end{itemize}
\end{columns}
\end{frame}

\begin{frame}[t]{Domain 3: Specialization Project}
The specialization requirement is a \textbf{semester-long Engineering Data Analytics project}.

\vspace{0.25cm}
Project options include:
\begin{itemize}[<+->]
  \tightitems
  \item Senior Design with a substantial data-analytics component
  \item Undergraduate research or independent research
  \item ENGR 2451 or ENGR 2453, when allowed by the certificate rules
  \item CEE 1323, IE 1171, IE 2064, or another approved project-based class
\end{itemize}

\vspace{0.2cm}
\uncover<5->{Students should communicate early with the instructor, senior-design advisor, or research advisor about the certificate requirement.}
\end{frame}

\begin{frame}[t]{What Counts as a Data Analytics Project?}
A qualifying project should involve writing code and performing a defensible data analysis, not only using spreadsheets or existing software.

\vspace{0.25cm}
\begin{itemize}[<+->]
  \tightitems
  \item Define the engineering or scientific problem.
  \item Describe the dataset, variables, observations, and data source.
  \item Clean and explore the data using plots and summaries.
  \item Select, justify, and interpret a statistical or predictive model.
  \item Challenge the results and discuss limitations.
  \item Write a report with methods, results, interpretation, and conclusions.
\end{itemize}
\end{frame}

\begin{frame}[t]{Certificate vs. IE Data Analytics Concentration}
\begin{center}
\begin{tabularx}{\textwidth}{@{}lYY@{}}
\toprule
 & \textbf{Certificate} & \textbf{IE Concentration}\\
\midrule
Who can pursue it? & Engineering students in any department and even non-engineering students & IE students only\\
Credits & 15 credits & 9 credits\\
Interdisciplinary rule & 9 credits must be beyond major requirements & More focused within IE\\
Project requirement & Includes a specialization project & Does not have the same certificate specialization structure\\
\bottomrule
\end{tabularx}
\end{center}
\end{frame}

\begin{frame}[t]{How to Plan Your Pathway}
\begin{itemize}[<+->]
  \smallitem
  \item Start with the published course list, then compare it with your major requirements.
  \item Make sure no more than two major courses are counted toward the certificate.
  \item If you want a new course to count, send the course number and description to Dr. Leu and Dr. Streiner.
  \item Course lists are expected to be updated periodically as Pitt adds new data-science offerings.
  \item For project-based courses, confirm the data-analytics component early.
\end{itemize}
\end{frame}

\begin{frame}[t]{How to Apply}
\begin{enumerate}[<+->]
  \smallitem
  \item Use the SSOE form for adding minors and certificates to your plan of study.
  \item Ask your undergraduate program coordinator or administrator for the form link.
  \item It is usually best to wait until your final semester before graduating to add the certificate to your plan.
  \item Check that the certificate appears correctly on your Academic Advising Report (AAR).
  \item If a course does not seem to count correctly, contact Dr. Leu or Dr. Streiner.
\end{enumerate}
\end{frame}

\begin{frame}[t]{ENGR 1451/2451: Exploratory Data Science}
\begin{columns}[T,onlytextwidth]
\column{0.50\textwidth}
\textbf{Counts toward the certificate}\par
\vspace{0.25cm}
Exploratory Data Science helps students turn real-world data into real-world answers.

\vspace{0.25cm}
\begin{itemize}[<+->]
  \tightitems
  \item Data exploration
  \item Visualization
  \item Predictive model development
  \item Statistical reasoning
\end{itemize}

\column{0.46\textwidth}
\begin{block}{Fall 2026}
Tuesday and Thursday\par
9:30--10:45 AM\par
Benedum 1047
\end{block}
\begin{block}{Questions}
Prof. Paul Leu\par
\href{mailto:pleu@pitt.edu}{pleu@pitt.edu}
\end{block}
\end{columns}
\end{frame}

\begin{frame}[t]{Where to Find More Information}
\begin{columns}[T,onlytextwidth]
\column{0.54\textwidth}
\textbf{Explore courses, requirements, and application details.}

\vspace{0.25cm}
\begin{itemize}[<+->]
  \item Website: Engineering \textrightarrow{} Industrial Engineering \textrightarrow{} Certificates
  \item Ask your undergraduate program coordinator for planning and form details.
  \item Contact Dr. Leu or Dr. Streiner for certificate-specific questions.
\end{itemize}

\vspace{0.2cm}
\href{https://www.engineering.pitt.edu/departments/industrial/undergraduate/degrees/certificates/}{Certificate website}

\column{0.40\textwidth}
\centering
\qrcode[height=3.3cm]{https://www.engineering.pitt.edu/departments/industrial/undergraduate/degrees/certificates/}

\vspace{0.2cm}
\small Scan for certificate information
\end{columns}
\end{frame}

\begin{frame}[t]{Questions?}
\centering
\Huge Engineering Data Analytics Certificate\\[0.35cm]
\Large Programming \quad Statistics \quad AI/ML\\[0.7cm]
\normalsize
Prof. Paul W. Leu\par
\href{mailto:pleu@pitt.edu}{pleu@pitt.edu}
\end{frame}

\end{document}
